\section{What exact WKB changed in the present research program}
\label{sec:program-development}
% BEGIN CURATED SUBJECT INDEX
\index{operator!adjoint and self-adjoint}
\index{spectrum!point}
\index{connection coefficient}
\index{boundary condition!incoming and outgoing}
\index{resonance!residue}
\index{Gamow--Siegert state}
\index{exact WKB}
\index{monodromy}
\index{Langer correction}
\index{Frobenius solution}
% END CURATED SUBJECT INDEX

The novelty of the program is not exhausted by obtaining several new pole
tables.  Exact WKB changed which pieces of the problem are regarded as
primary.  Before choosing a Hilbert-space realization one can construct an
analytic incoming coefficient, follow it across Stokes sectors, and ask which
features survive every admissible representation.  This produces the
following shifts of emphasis:
\begin{longtable}{@{}p{0.29\linewidth}p{0.63\linewidth}@{}}
\toprule
conventional starting point & connection-first understanding \\
\midrule
\endfirsthead
\toprule
conventional starting point & connection-first understanding \\
\midrule
\endhead
complex eigenvalue & zero of a globally transported incoming coefficient;
the eigenvalue is one realization of the zero \\
\addlinespace
divergent Gamow norm & derivative or residue of the analytic connection datum;
regularizations must reproduce it \\
\addlinespace
isolated resonance level & pole and continuum as parts of one continued
spectral resolution \\
\addlinespace
Langer correction as a prescription & local singular monodromy required by
global consistency \\
\addlinespace
EP square root of a matrix & multiple zero of a scattering determinant with
continuum normalization kept explicit \\
\addlinespace
black-hole damping boundary condition & horizon-to-horizon connection problem
on a complex spectral curve \\
\bottomrule
\end{longtable}

These shifts have several consequences for how the theory is used.

First, a resonance can be constructed before choosing a non-self-adjoint
operator.  Exact WKB transports local solutions, produces
$\Cincoming(E)$, and defines a simple pole by
$\Cincoming(\resonantE)=0$ together with
$\Cincoming'(\resonantE)\neq0$.  A complex-scaled eigenvalue is then
a calculational realization of this datum.  This order prevents an arbitrary
complex eigenvalue of a truncation from being mistaken for a resonance.

Second, equivalence of regularizations becomes a checkable residue identity.
The Zel'dovich pairing, the rigged-space functional, and the left--right
complex-scaled dyad must reproduce the derivative of the same connection
coefficient.  Agreement of pole positions without agreement of the residue is
only a partial test.

Third, the continuum is not a background that may be deleted after a pole is
found.  Rotating the continuum, deforming a Berggren contour, or extracting a
resonance term reorganizes one spectral response.  Threshold cuts, continuum
level density, and source-dependent excitation survive that reorganization
and control where a pole approximation is accurate.

Fourth, new applications are classified by the connection data they add.  A
radial origin adds Frobenius indices and singular monodromy; an exceptional
point changes a simple zero into a multiple zero and requires a Jordan chain;
a black-hole horizon replaces spatial infinity by a second radiation sector.
The local differential equation may still look Schr\"odinger-like, but the
global connection problem has changed and must be rebuilt accordingly.

Finally, the same organization separates theorem, representation, and
numerics.  Analytic dilation supplies the theorem-controlled wedge.  A basis
and scaling angle supply a representation.  Stability under the parameters
that should be inessential supplies a numerical falsification test.  The
three layers cooperate, but none can be substituted silently for another.

The research articles in which these calculations were first developed are
Refs.~\cite{Morikawa:2025grx,Morikawa:2025xjq,Morikawa:2025vvs,
Morikawa:2025ezl,Morikawa:2025inx,Ogawa:2026veu,Ogawa:2026dSCSM}.
Their chronology is bibliographic provenance; the logical order used here is
the dependency order of the connection problem.
